%============================================================
\section{Computational continuation and Floquet diagnostics}
\label{sec:numerics}
%============================================================

This section examines the three local branches obtained in
Sections~\ref{sec:local-continuation}--\ref{sec:stability-quantitative}
by numerical continuation from the exact constant states at zero
forcing. The computations are used to compare the periodic responses
with the asymptotic expansions, to monitor their Floquet spectra, and
to assess discretization sensitivity. They do not constitute a
computer-assisted existence proof.

All computations were performed in Python using \texttt{NumPy},
\texttt{SciPy}, and \texttt{Matplotlib}; see
\cite{VirtanenEtAl2020}. The initial-value problems were integrated with
the explicit DOP853 Runge--Kutta method through
\texttt{scipy.integrate.solve\_ivp}, with relative and absolute
tolerances
\[
\texttt{rtol}=10^{-10},
\qquad
\texttt{atol}=10^{-12}.
\]
The implementation, input parameters, generated tables, and figures are
provided in the companion repository
\cite{LopezCastroRamirezCode2026}.

%------------------------------------------------------------
\subsection{Shooting and parameter continuation}
\label{subsec:numerical-shooting}
%------------------------------------------------------------

For the harmonic problem
\begin{equation}
\label{eq:numerical-harmonic-problem}
 u''+\delta u'+f_{\alpha}(u)=\beta\cos(\omega t),
 \qquad
 T=\frac{2\pi}{\omega},
\end{equation}
we write
\[
 x=(u,v)^{\mathsf T}
\]
and denote by
\[
 P_{\beta}(x_0)
\]
the time-$T$ map of the associated first-order system. A periodic
solution corresponds to a zero of
\begin{equation}
\label{eq:numerical-shooting-residual}
 G_{\beta}(x_0)
 :=
 P_{\beta}(x_0)-x_0.
\end{equation}
The derivative of the shooting residual is
\begin{equation}
\label{eq:numerical-shooting-jacobian}
 D G_{\beta}(x_0)
 =
 M_{\beta}(x_0)-I_2,
\end{equation}
where $M_{\beta}(x_0)$ is obtained by integrating the variational
equation simultaneously with the orbit. A damped Newton correction is
then applied to \eqref{eq:numerical-shooting-residual}.

At $\beta=0$, the three exact initial states are
\begin{equation}
\label{eq:numerical-initial-states}
 x_-^{0}=(-\sqrt{\alpha^2-1},0),
 \qquad
 x_0^{0}=(0,0),
 \qquad
 x_+^{0}=(\sqrt{\alpha^2-1},0).
\end{equation}
Starting from these states, each branch is continued on a uniform
parameter grid
\[
 \beta_n=n\,\Delta\beta,
 \qquad
 \Delta\beta=10^{-2},
\]
using the solution at $\beta_n$ as the predictor for the Newton
correction at $\beta_{n+1}$. This continuation procedure is appropriate
as long as $\beta$ remains a regular local parameter. It is not designed
to cross folds; a pseudo-arclength formulation would be required near a
singular shooting Jacobian; see \cite{AllgowerGeorg1990,Seydel2010}.

The numerical example uses
\begin{equation}
\label{eq:numerical-parameters}
 \alpha=2,
 \qquad
 \delta=0.35,
 \qquad
 \omega=1,
 \qquad
 T=2\pi.
\end{equation}
The three branches were continued over
\begin{equation}
\label{eq:numerical-beta-range}
 0\leq\beta\leq0.40.
\end{equation}
This interval is a computational continuation range. It is not asserted
to be an analytic persistence interval and is not identified with the
geometric quantity $M_{\alpha}$.

Figure~\ref{fig:numerical-branches} displays the minimum and maximum of
each computed solution over one period. The continuation preserves three
disjoint response families throughout the interval
\eqref{eq:numerical-beta-range}.

\begin{figure}[htbp]
\centering
\includegraphics[width=0.88\linewidth]
{figures/branch_continuation.pdf}
\caption{Minimum and maximum over one period along the three branches
continued from the constant states at $\beta=0$.}
\label{fig:numerical-branches}
\end{figure}

%------------------------------------------------------------
\subsection{Comparison with the local asymptotic expansions}
\label{subsec:numerical-asymptotics}
%------------------------------------------------------------

Let $z_j$ and $w_j$ be the first- and second-order terms in
Theorem~\ref{thm:harmonic-second-order-expansion}, and let $\chi_0$ be
the cubic correction in
Theorem~\ref{thm:central-cubic-expansion}. We define
\begin{align}
 E_{+,1}(\beta)
 &:={}
 \left\|
 u_+(\cdot;\beta)
 -\xi_+
 -\beta z_+
 \right\|_{\infty},
 \label{eq:numerical-error-plus-one}
 \\
 E_{+,2}(\beta)
 &:={}
 \left\|
 u_+(\cdot;\beta)
 -\xi_+
 -\beta z_+
 -\beta^2w_+
 \right\|_{\infty},
 \label{eq:numerical-error-plus-two}
 \\
 E_{0,1}(\beta)
 &:={}
 \left\|
 u_0(\cdot;\beta)
 -\beta z_0
 \right\|_{\infty},
 \label{eq:numerical-error-zero-one}
 \\
 E_{0,3}(\beta)
 &:={}
 \left\|
 u_0(\cdot;\beta)
 -\beta z_0
 -\beta^3\chi_0
 \right\|_{\infty}.
 \label{eq:numerical-error-zero-three}
\end{align}
The negative exterior branch gives the same errors, up to the numerical
integration tolerance, by the symmetry established in
Corollary~\ref{cor:smooth-harmonic-branches}.

\begin{table}[htbp]
\centering
\caption{Uniform errors of the asymptotic approximations.}
\label{tab:numerical-asymptotic-errors}
\small
\setlength{\tabcolsep}{8pt}
\begin{tabular}{ccccc}
\hline
$\beta$
& $E_{+,1}(\beta)$
& $E_{+,2}(\beta)$
& $E_{0,1}(\beta)$
& $E_{0,3}(\beta)$
\\
\hline
$0.02$
& $3.0402\times10^{-4}$
& $2.2270\times10^{-5}$
& $3.7554\times10^{-7}$
& $4.6950\times10^{-11}$
\\
$0.05$
& $2.0709\times10^{-3}$
& $3.4984\times10^{-4}$
& $5.8695\times10^{-6}$
& $1.7461\times10^{-9}$
\\
$0.10$
& $9.4967\times10^{-3}$
& $2.8190\times10^{-3}$
& $4.6997\times10^{-5}$
& $5.5717\times10^{-8}$
\\
\hline
\end{tabular}
\end{table}

The values in Table~\ref{tab:numerical-asymptotic-errors} are consistent
with
\[
 E_{+,1}(\beta)=O(\beta^2),
 \qquad
 E_{+,2}(\beta)=O(\beta^3),
\]
and
\[
 E_{0,1}(\beta)=O(\beta^3),
 \qquad
 E_{0,3}(\beta)=O(\beta^5).
\]
The absence of quadratic and quartic corrections on the central branch
is therefore visible at the numerical level.

%------------------------------------------------------------
\subsection{Representative profiles and Floquet spectra}
\label{subsec:numerical-representative-case}
%------------------------------------------------------------

We now fix
\begin{equation}
\label{eq:numerical-representative-beta}
 \beta=0.20.
\end{equation}
The three periodic profiles obtained by continuation are shown in
Figure~\ref{fig:numerical-profiles}. Their initial data, ranges, and
Floquet multipliers are reported in
Table~\ref{tab:numerical-floquet-data}.

\begin{figure}[htbp]
\centering
\includegraphics[width=0.88\linewidth]
{figures/representative_orbits.pdf}
\caption{Three periodic responses at $\beta=0.20$ for the parameter set
\eqref{eq:numerical-parameters}.}
\label{fig:numerical-profiles}
\end{figure}

\begin{table}[htbp]
\centering
\caption{Shooting and Floquet data at $\beta=0.20$.}
\label{tab:numerical-floquet-data}
\small
\setlength{\tabcolsep}{5pt}
\begin{tabular}{lcccl}
\hline
Branch
& $(u(0),u'(0))$
& $[\min u,\max u]$
& Floquet multipliers
& Classification
\\
\hline
$u_-$
& $(-1.982454,\,0.366600)$
& $[-2.164979,-1.262030]$
& $0.146329\pm0.299147\,\mathrm{i}$
& exponentially stable
\\
$u_0$
& $(-0.097381,\,0.017143)$
& $[-0.098876,0.098876]$
& $187.535111,\;5.91363\times10^{-4}$
& hyperbolically unstable
\\
$u_+$
& $(1.429238,\,0.347437)$
& $[1.262030,2.164979]$
& $0.146329\pm0.299147\,\mathrm{i}$
& exponentially stable
\\
\hline
\end{tabular}
\end{table}

The modulus of each exterior multiplier is
\[
 0.333018<1,
\]
whereas the central branch has one multiplier greater than one and one
smaller than one. This agrees with
Theorem~\ref{thm:stability-persistence}. The exterior spectra coincide,
as predicted by Proposition~\ref{prop:exterior-floquet-symmetry}.

For all three solutions,
\[
 e^{-\delta T}
 =
 e^{-0.35(2\pi)}
 =
 0.110901278364\ldots.
\]
The computed determinants of the monodromy matrices differ from this
value by
\[
 3.66\times10^{-12},
 \qquad
 1.24\times10^{-11},
 \qquad
 3.86\times10^{-12}
\]
for $u_-$, $u_0$, and $u_+$, respectively. The corresponding shooting
residuals are
\[
 2.00\times10^{-15},
 \qquad
 1.89\times10^{-13},
 \qquad
 2.03\times10^{-12}.
\]

The half-period identities were also evaluated on a uniform grid of
$4096$ points. The resulting defects were
\begin{equation}
\label{eq:numerical-symmetry-defects}
 \left\|
 u_-(t;\beta)
 +u_+\left(t+\frac{T}{2};\beta\right)
 \right\|_{\infty}
 =1.50\times10^{-9}
\end{equation}
and
\begin{equation}
\label{eq:numerical-central-symmetry-defect}
 \left\|
 u_0(t;\beta)
 +u_0\left(t+\frac{T}{2};\beta\right)
 \right\|_{\infty}
 =8.51\times10^{-11}.
\end{equation}

Figure~\ref{fig:numerical-floquet-continuation} shows the spectral radius
of the three monodromy matrices along the computed continuation
interval. The exterior branches remain inside the unit disk, while the
central branch remains hyperbolically unstable over the sampled
parameter range.

\begin{figure}[htbp]
\centering
\includegraphics[width=0.88\linewidth]
{figures/floquet_moduli.pdf}
\caption{Spectral radius of the monodromy matrix along the computed
branches. The dashed line marks the unit-modulus threshold.}
\label{fig:numerical-floquet-continuation}
\end{figure}

%------------------------------------------------------------
\subsection{Independent finite-difference check}
\label{subsec:numerical-finite-difference}
%------------------------------------------------------------

As an independent discretization, the periodic boundary value problem
was approximated on the uniform grid
\[
 t_k=\frac{kT}{N},
 \qquad
 k=0,\ldots,N-1,
\]
using centered periodic differences,
\begin{equation}
\label{eq:numerical-centered-differences}
 (D_1U)_k
 =
 \frac{U_{k+1}-U_{k-1}}{2\Delta t},
 \qquad
 (D_2U)_k
 =
 \frac{U_{k+1}-2U_k+U_{k-1}}{\Delta t^2},
\end{equation}
with indices interpreted modulo $N$. The nonlinear system
\begin{equation}
\label{eq:numerical-discrete-system}
 D_2U
 +\delta D_1U
 +f_{\alpha}(U)
 -\beta H
 =0,
 \qquad
 H_k=\cos(\omega t_k),
\end{equation}
was solved by sparse Newton iteration with the exact discrete Jacobian.
The shooting profiles were used only as initial approximations for this
independent algebraic solve; see \cite{AscherMattheijRussell1995}.

\begin{table}[htbp]
\centering
\caption{Amplitude intervals from the centered finite-difference
scheme at $\beta=0.20$.}
\label{tab:numerical-mesh-intervals}
\small
\begin{tabular}{cccc}
\hline
$N$ & $u_N^-$ & $u_N^0$ & $u_N^+$
\\
\hline
$128$
& $[-2.165131,-1.261909]$
& $[-0.098862,0.098862]$
& $[1.261909,2.165131]$
\\
$256$
& $[-2.165013,-1.261971]$
& $[-0.098879,0.098879]$
& $[1.261971,2.165013]$
\\
$512$
& $[-2.164989,-1.262016]$
& $[-0.098877,0.098877]$
& $[1.262016,2.164989]$
\\
\hline
\end{tabular}
\end{table}

For nested grids, define
\begin{equation}
\label{eq:numerical-nested-grid-error}
 \Delta_{N,2N}(U)
 :=
 \max_{0\leq k<N}
 |U_N(t_k)-U_{2N}(t_{2k})|.
\end{equation}
The computed values are
\begin{table}[htbp]
\centering
\caption{Nested-grid discrepancies for the three branches.}
\label{tab:numerical-mesh-discrepancies}
\small
\begin{tabular}{lccc}
\hline
Transition
& $\Delta_{N,2N}(u_-)$
& $\Delta_{N,2N}(u_0)$
& $\Delta_{N,2N}(u_+)$
\\
\hline
$128\to256$
& $1.8486\times10^{-4}$
& $9.0429\times10^{-6}$
& $1.8486\times10^{-4}$
\\
$256\to512$
& $4.6206\times10^{-5}$
& $2.2607\times10^{-6}$
& $4.6206\times10^{-5}$
\\
\hline
\end{tabular}
\end{table}

The discrepancies decrease by a factor close to four after mesh
doubling, consistently with the second-order centered discretization.
The discrete residuals remained below $8.6\times10^{-11}$ for all
reported meshes and branches.

%------------------------------------------------------------
\subsection{Scope of the computations}
\label{subsec:numerical-scope}
%------------------------------------------------------------

The computations establish the following numerical facts for the
parameter set \eqref{eq:numerical-parameters}: the three branches can be
continued from $\beta=0$ to $\beta=0.40$ by regular parameter
continuation; the local asymptotic expansions display the predicted
orders; the exterior branches remain exponentially stable and share the
same Floquet spectrum; the central branch remains hyperbolically
unstable; and the periodic profiles are stable under mesh refinement.

These conclusions concern the computed branches and parameter grid. No
claim is made that $\beta=0.40$ is a maximal continuation value, a
bifurcation threshold, or an analytically certified coexistence bound.
The residual, Liouville, symmetry, and refinement tests are
reproducibility diagnostics rather than interval-certified existence
proofs.

\FloatBarrier
